%
% Capítulo 1
%
\chapter{Introdução} \label{cap:intro}

Este capítulo tem como objetivo apresentar a motivação do projeto e as necessidades que este irá preencher.

%
% Secção 1.1
%
\section{Motivação} \label{sec11}

Durante a fase de definição do projeto, surgiu uma proposta por parte do presidente do Grupo Desportivo União Ericeirense, dirigida aos autores deste trabalho, então a concluir a Licenciatura em Engenharia Informática e de Computadores. Durante a apresentação da proposta, tornou-se evidente a existência de uma necessidade concreta: reforçar a ligação entre o clube e os adeptos.

A ausência de uma plataforma digital representava, para o clube, a perda de diversas oportunidades. Muitos potenciais sócios e patrocinadores acabavam por não concretizar a sua adesão devido à reduzida disponibilidade para se deslocarem presencialmente à sede do clube. Identificou-se, assim, uma oportunidade clara de intervenção, capaz de colmatar estas dificuldades.

A solução inicialmente equacionada passava pelo desenvolvimento de uma aplicação móvel. Contudo, a burocracia e os requisitos de licenciamento associados à publicação em plataformas como a App Store e a Google Play levaram à adoção de uma alternativa mais acessível: uma aplicação web. Esta abordagem permite disponibilizar as funcionalidades pretendidas sem as restrições inerentes à distribuição em lojas de aplicações móveis.

%
% Secção 1.2
%
\section{Objetivos} \label{sec12}

O objetivo central deste projeto é desenvolver uma web app que reforce a ligação entre o Grupo Desportivo União Ericeirense e a sua comunidade, disponibilizando online um conjunto de processos que tradicionalmente exigiam deslocação presencial à sede do clube.

Para concretizar este propósito, definiram-se os 4 objetivos, sendo o último o que os dirigentes do clube consideraram menos crucial, no entanto enriquecia a complexidade do projeto:

\begin{itemize}
    \item Permitir que qualquer pessoa se torne sócio do clube através da aplicação web;
    \item Permitir a adesão de novos patrocinadores de forma totalmente online;
    \item Disponibilizar a inscrição de atletas no clube;
    \item Implementar a venda de bilhetes para os jogos, com validação por código QR.
\end{itemize}

%
% Secção 1.3
%
\section{Terminologia} \label{sec13}

Esta secção é dedicada à definição de termos usados ao longo deste documento.

\begin{itemize}
    \item \textit{frontend} -- interface, parte da aplicação que é executada no navegador do utilizador;
    \item \textit{backend} -- componente da aplicação que é executada no servidor;
    \item \textit{backoffice} -- área de gestão interna utilizada pelo clube;
    \item \textit{role} -- nível de permissão atribuído a um utilizador;
    \item \textit{check-in} -- validação da entrada através da leitura do código QR;
    \item \textit id -- identificador único;
    \item \textit QR Code -- \textit{Quick Response Code}, código de barras bidimensional;
\end{itemize}